\section{Introduction}
\label{sec:intro}

Silicon photonic circuits are designed around nominal waveguide cross-sections, and their behavior follows the
fabricated ones. Width and thickness variations across a wafer shift ring resonances and interferometer phases, and
models of layout-dependent variability are now used to predict circuit yield~\cite{lu2017}. The sidewall angle is
the least monitored of these parameters. Plasma etching leaves walls that lean away from vertical, the lean varies
with etch chemistry, pattern density and position on the wafer \todo{cite: measured sidewall-angle variation in
silicon photonic etch processes}, and it moves the effective index of every mode that touches the wall. Circuits
that mix fully etched strips with partially etched ribs see two different angles from two different etch steps, so
a component designed for one cross-section can sit off resonance when either angle drifts.

Dimensional metrology measures the angle directly but sparsely. Cross-sectional electron microscopy requires cutting
the sample and is destructive~\cite{orji2018}; atomic force microscopy needs specialized tips and scans to reach a
wall \todo{cite: 3D/CD-AFM sidewall metrology on Si waveguides}; optical scatterometry is fast and non-destructive but infers the profile through a parametrized model whose
parameters can be correlated and whose inverse problem need not have a unique solution~\cite{orji2018,raymond2005}.
All of these measure shape at a few sites. What a circuit designer needs is the optical consequence of the shape
across the wafer, which suggests measuring the angle with the same ring and interferometer test structures that
already monitor width and thickness. Such optical extraction methods recover width and thickness from measured
effective and group indices~\cite{xing2018,dwivedi2015}, but they hold the sidewall angle at an assumed value
\todo{verify in the full text of Xing et al.\ 2018 that the extraction fixes the sidewall angle (search snippet: 85$^\circ$ from a foundry handbook)}.

The reason the angle is usually fixed is not a lack of data but a property of the structures. This paper shows that
for oxide-clad strip waveguides the angle is not merely hard to separate from the width: to first order it cannot be
separated at all, for any number of strips, polarizations or readouts. We then ask which structures restore the
missing information, how much prior knowledge of the process they need, and what precision they reach under a
measurement noise taken from published ring-resonator metrology~\cite{janz2023}. The answers come from full-vector
finite-element mode solutions and from the Cram\'er--Rao lower bound (CRLB) of a joint fit, checked by nonlinear
fits of simulated noisy data. No devices were fabricated; every number below is a simulation result or an input
taken from the cited literature.

The contributions are:
\begin{itemize}
\item a first-order proof, confirmed numerically, that oxide-clad strips are exactly blind to the sidewall angle,
with the angle/width ratio $R=(y_c-y_\mathrm{pivot})/(h/2)$ set by the wall-perturbation centroid of the mode;
\item an identifiability analysis showing that one rib made by a separate etch cannot determine the strip angle,
and that prior knowledge of the etch depth does not change this;
\item test-structure sets built from several widths of the same rib that reach \qty{0.39}{\degree} (two rib widths)
and \qtyrange{0.12}{0.14}{\degree} (best set of eight) at a \qty{1}{pm} wavelength precision, with the bound confirmed
by Monte Carlo fits;
\item a methods comparison showing that finite-difference derivatives from a grid-based solver were not accurate
enough for this analysis, and that small derivative errors in any solver create false angle information unless the
strip symmetry is imposed.
\end{itemize}
